\section{Trabajos relacionados}
La arquitectura U-Net \cite{ronneberger2015unet} es el punto de partida dominante en segmentación biomédica por su combinación de encoder-decoder y conexiones de salto que preservan detalle espacial. Variantes posteriores como Attention U-Net \cite{oktay2018attention} y U-Net++ \cite{zhou2019unetpp} mejoran la fusión multiescala y el filtrado de información irrelevante en las skip connections.

En tareas de segmentación vascular, la literatura reporta que la elección de la función de pérdida es determinante debido al fuerte desbalance fondo/vaso. Dice loss \cite{milletari2016vnet} optimiza directamente superposición, BCE estabiliza el entrenamiento y Focal loss \cite{lin2017focal} enfatiza ejemplos difíciles, particularmente vasos finos de bajo contraste.

Sobre los datasets, DRIVE \cite{staal2004drive} y CHASE\_DB1 son referentes para evaluar tanto rendimiento in-distribution como capacidad de generalización. Revisiones amplias del campo \cite{fraz2012retinal} muestran que las caídas entre datasets siguen siendo una limitación abierta, motivando técnicas de adaptación y preprocesamiento. En esa línea, métodos de realce local como CLAHE suelen emplearse para homogenizar contraste y mejorar la detectabilidad de capilares.
